\documentclass[11pt,a4paper]{article}

\usepackage[utf8]{inputenc}
\usepackage[T1]{fontenc}
\usepackage[french]{babel}
\usepackage{geometry}
\usepackage{booktabs}
\usepackage{longtable}
\usepackage{array}
\usepackage{xcolor}
\usepackage{colortbl}
\usepackage{graphicx}
\usepackage{fancyhdr}
\usepackage{titlesec}
\usepackage{tikz}
\usepackage{enumitem}
\usepackage{hyperref}
\usepackage{multicol}
\usepackage{listings}
\usepackage{tcolorbox}
\usepackage{amssymb}
\usepackage{float}
\usepackage{caption}
\usepackage{setspace}
\usepackage{amsmath}
\usepackage{multirow}

\geometry{margin=2cm}
\pagestyle{fancy}
\fancyhf{}
\rhead{Atelier Machine Learning}
\lhead{Analyse Comportementale Clientèle}
\rfoot{Page \thepage}

\definecolor{headerblue}{RGB}{41, 65, 114}
\definecolor{lightgray}{RGB}{245, 245, 245}
\definecolor{accentorange}{RGB}{230, 126, 34}
\definecolor{codegreen}{RGB}{0,128,0}
\definecolor{codegray}{RGB}{128,128,128}
\definecolor{codepurple}{RGB}{128,0,128}
\definecolor{backcolour}{RGB}{248,248,248}
\definecolor{warningred}{RGB}{231, 76, 60}
\definecolor{successgreen}{RGB}{46, 204, 113}

\lstdefinestyle{mystyle}{
    backgroundcolor=\color{backcolour},
    commentstyle=\color{codegreen},
    keywordstyle=\color{codepurple},
    numberstyle=\tiny\color{codegray},
    stringstyle=\color{codegreen},
    basicstyle=\ttfamily\footnotesize,
    breakatwhitespace=false,
    breaklines=true,
    captionpos=b,
    keepspaces=true,
    numbers=left,
    numbersep=5pt,
    showspaces=false,
    showstringspaces=false,
    showtabs=false,
    tabsize=2
}
\lstset{style=mystyle}

\tcbset{
    colback=backcolour,
    colframe=headerblue,
    fonttitle=\bfseries,
    boxrule=1pt,
    arc=3mm
}

\titleformat{\section}
{\color{headerblue}\normalfont\Large\bfseries}
{\thesection}{1em}{}[{\titlerule[2pt]}]

\titleformat{\subsection}
{\color{accentorange}\normalfont\large\bfseries}
{\thesubsection}{1em}{}

\newcolumntype{C}[1]{>{\centering\arraybackslash}p{#1}}
\newcolumntype{L}[1]{>{\raggedright\arraybackslash}p{#1}}

\setstretch{1.15}

\begin{document}

% ================= PAGE DE GARDE =================
\begin{titlepage}
\begin{tikzpicture}[remember picture,overlay]
    \fill[headerblue] (current page.north west) rectangle ([yshift=-5cm]current page.north east);
    \node[anchor=north west, xshift=2cm, yshift=-2cm, text=white] at (current page.north west)
        {\Huge\textbf{Atelier Machine Learning}};
    \node[anchor=north west, xshift=2cm, yshift=-3.5cm, text=white] at (current page.north west)
        {\Large\textit{Analyse Comportementale de la Clientèle Retail}};
\end{tikzpicture}

\vspace{6cm}

\begin{center}
{\LARGE\textbf{COMPTE RENDU FINAL}}\par
\vspace{1cm}
{\large Atelier Pratique : Prédiction du churn client}\par
\vspace{2cm}
{\large \textbf{Préparé par :}}\par
\vspace{0.4cm}
{\Large \textbf{Yassine Meftah}}

\vspace{3cm}

\begin{tikzpicture}
    \draw[headerblue, line width=2pt] (0,0) -- (10,0);
\end{tikzpicture}

\vspace{2cm}

{\textbf{Année Universitaire : 2025-2026}}
\end{center}
\vfill
\end{titlepage}

% ================= RESUME =================
\section*{Résumé}
Ce projet s’inscrit dans le cadre de l’analyse comportementale de la clientèle d’une plateforme de e-commerce. L’objectif principal est de construire une chaîne complète de Machine Learning capable de prédire le churn client, d’explorer la structure des données, de segmenter les profils clients et de proposer une solution d’inférence exploitable. Un pipeline complet a été mis en place : préparation des données, parsing des dates, feature engineering, détection et correction des valeurs aberrantes, suppression des colonnes redondantes ou fuyantes, rééquilibrage des classes par SMOTENC, réduction de dimension par ACP, clustering KMeans, classification supervisée, régression et validation croisée. Le meilleur modèle de classification obtenu est un \texttt{RandomForest\_GridSearch}. Le meilleur modèle de régression obtenu est un \texttt{HistGradientBoostingRegressor}. Une phase d’inférence a également été réalisée via un script de prédiction et une interface Flask. Le projet met ainsi en œuvre une démarche complète, reproductible et proche d’un cas réel de mise en production.

\section*{Abstract}
This project focuses on customer behavioral analysis in a retail e-commerce setting. The main objective is to build a complete Machine Learning pipeline capable of predicting customer churn, exploring the structure of the data, segmenting customer profiles, and providing a reusable inference solution. A full workflow was implemented: data preparation, date parsing, feature engineering, outlier handling, removal of redundant and leakage-prone variables, class rebalancing with SMOTENC, dimensionality reduction with PCA, KMeans clustering, supervised classification, regression, and cross-validation. The best classification model is a \texttt{RandomForest\_GridSearch}, while the best regression model is a \texttt{HistGradientBoostingRegressor}. An inference phase was also developed through a prediction script and a Flask interface. Overall, this project demonstrates a complete, reproducible workflow close to a real deployment scenario.

\tableofcontents
\newpage

% ================= INTRODUCTION =================
\section{Introduction générale}

\subsection{Contexte}
Dans le domaine du retail et du e-commerce, l’analyse des données clients constitue un levier stratégique pour améliorer la fidélisation, personnaliser les campagnes marketing et optimiser la rentabilité. La capacité à identifier les clients susceptibles de quitter la plateforme permet d’anticiper les actions commerciales et de réduire la perte de valeur.

\subsection{Problématique}
Le churn client représente un enjeu majeur pour les entreprises, car il est généralement plus coûteux d’acquérir un nouveau client que de conserver un client existant. Le problème traité dans ce projet consiste donc à exploiter les données comportementales et transactionnelles afin de prédire si un client présente un risque de départ.

\subsection{Objectifs du projet}
Le projet poursuit plusieurs objectifs complémentaires :
\begin{itemize}
    \item développer un modèle de classification pour prédire le churn ;
    \item analyser la structure des données et identifier les variables importantes ;
    \item appliquer une réduction de dimension par ACP pour la visualisation et le clustering ;
    \item segmenter les clients via KMeans ;
    \item entraîner un modèle de régression pour prédire \texttt{MonetaryTotal} ;
    \item déployer une solution d’inférence réutilisable via un script Python et une interface Flask.
\end{itemize}

\subsection{Démarche suivie}
La méthodologie suivie s’organise selon une chaîne complète :
\begin{center}
\textbf{Préparation des données $\rightarrow$ Transformation $\rightarrow$ Modélisation $\rightarrow$ Évaluation $\rightarrow$ Déploiement}
\end{center}

% ================= DONNEES =================
\section{Présentation du jeu de données}

\subsection{Description générale}
Le dataset étudié correspond à une base client retail contenant des variables numériques, catégorielles et dérivées. Il inclut des informations transactionnelles, temporelles, démographiques et comportementales. La variable cible principale du projet est \texttt{Churn}.

\subsection{Typologie des variables}
Les variables peuvent être regroupées en plusieurs familles :
\begin{itemize}
    \item \textbf{Variables transactionnelles :} \texttt{Frequency}, \texttt{MonetaryTotal}, \texttt{MonetaryAvg}, \texttt{MonetaryStd}, \texttt{UniqueProducts}, \texttt{ReturnRatio}, etc.
    \item \textbf{Variables temporelles :} \texttt{PreferredHour}, \texttt{PreferredMonth}, \texttt{RegYear}, \texttt{RegMonth}, \texttt{RegDay}, etc.
    \item \textbf{Variables de relation client :} \texttt{SupportTicketsCount}, \texttt{SatisfactionScore}.
    \item \textbf{Variables démographiques et contextuelles :} \texttt{Age}, \texttt{Gender}, \texttt{Country}, \texttt{Region}.
\end{itemize}

\subsection{Problèmes de qualité identifiés}
L’analyse initiale du jeu de données a mis en évidence :
\begin{itemize}
    \item des valeurs manquantes, notamment sur \texttt{Age} ;
    \item un format de date à normaliser pour \texttt{RegistrationDate} ;
    \item la présence de valeurs aberrantes sur certaines colonnes ;
    \item des variables redondantes ;
    \item des variables potentiellement fuyantes par rapport à la cible ;
    \item un déséquilibre entre les classes \texttt{Churn = 0} et \texttt{Churn = 1}.
\end{itemize}

% ================= EXPLORATION =================
\section{Analyse exploratoire et compréhension des données}

\subsection{Qualité des données}
La première étape a consisté à examiner la qualité globale du jeu de données, notamment la présence de valeurs manquantes, de colonnes peu informatives, de formats incohérents et de variables susceptibles de dégrader l’apprentissage.

\subsection{Corrélations et multicolinéarité}
Une analyse des corrélations entre variables numériques a été réalisée afin d’identifier les dépendances fortes et les redondances.

\begin{figure}[H]
    \centering
    \includegraphics[width=0.92\textwidth]{reports/correlation_heatmap_numeric.png}
    \caption{Matrice de corrélation des variables numériques après nettoyage.}
\end{figure}

Cette heatmap montre une structure plus saine après nettoyage : certaines corrélations modérées demeurent, ce qui est normal, mais les redondances les plus problématiques ont été réduites.

\subsection{Lecture métier}
Plusieurs variables ont une signification métier directe :
\begin{itemize}
    \item \texttt{Frequency} mesure l’intensité d’activité d’un client ;
    \item \texttt{MonetaryTotal} traduit la valeur économique générée ;
    \item \texttt{UniqueProducts} et \texttt{AvgProductsPerTransaction} décrivent la diversité et la richesse du panier ;
    \item \texttt{SupportTicketsCount} et \texttt{SatisfactionScore} renseignent sur la qualité de la relation client.
\end{itemize}

% ================= PREPARATION =================
\section{Préparation des données}

\subsection{Nettoyage initial}
Plusieurs colonnes ont été supprimées dès le départ :
\begin{itemize}
    \item les identifiants non prédictifs ;
    \item les variables constantes ;
    \item les variables fuyantes ou trop proches de la cible ;
    \item certaines variables redondantes détectées lors de l’analyse exploratoire.
\end{itemize}

\subsection{Détection de fuite d’information}
Une étape spécifique d’audit a été menée pour identifier les colonnes suspectes vis-à-vis de la cible, telles que :
\begin{itemize}
    \item \texttt{AccountStatus}
    \item \texttt{LoyaltyLevel}
    \item \texttt{SpendingCategory}
    \item \texttt{BasketSizeCategory}
    \item \texttt{AvgBasketValue}
\end{itemize}

L’objectif était de vérifier si certaines colonnes reconstruisaient implicitement la cible ou en contenaient une approximation trop forte.

\subsection{Parsing et transformation des dates}
La variable \texttt{RegistrationDate} a été convertie puis décomposée en plusieurs composantes :
\begin{itemize}
    \item \texttt{RegYear}
    \item \texttt{RegMonth}
    \item \texttt{RegDay}
    \item \texttt{RegWeekday}
\end{itemize}

\subsection{Feature engineering}
Un travail de feature engineering a été réalisé sur :
\begin{itemize}
    \item les adresses IP : création de \texttt{IsPrivateIP} et \texttt{IpFirstOctet} ;
    \item les variables comportementales existantes ;
    \item la structuration finale d’un ensemble de variables robuste et interprétable.
\end{itemize}

\subsection{Suppression de variables redondantes}
Des variables redondantes ont été supprimées afin de simplifier l’espace de représentation et limiter la duplication d’information. Cette étape repose à la fois sur l’analyse statistique et sur une lecture métier.

\subsection{Gestion des outliers}
Les valeurs aberrantes ont été traitées via la méthode IQR. Les bornes ont été calculées sur le jeu d’entraînement uniquement, puis réappliquées au jeu de test afin d’éviter toute fuite d’information.

\subsection{Gestion des valeurs manquantes}
Le pipeline de prétraitement combine plusieurs stratégies :
\begin{itemize}
    \item médiane pour la plupart des variables numériques ;
    \item mode pour les variables catégorielles ;
    \item \texttt{KNNImputer} spécifiquement pour \texttt{Age}.
\end{itemize}

\subsection{Encodage et normalisation}
Le prétraitement final repose sur un \texttt{ColumnTransformer} combinant :
\begin{itemize}
    \item \texttt{SimpleImputer} + \texttt{StandardScaler} pour les variables numériques ;
    \item \texttt{KNNImputer} + \texttt{StandardScaler} pour \texttt{Age} ;
    \item \texttt{SimpleImputer} + \texttt{OneHotEncoder} pour les variables catégorielles.
\end{itemize}

\subsection{Rééquilibrage des classes}
Le jeu d’entraînement a été rééquilibré par \texttt{SMOTENC}. Cette opération n’a été appliquée que sur le train, afin de respecter une méthodologie correcte.

\subsection{Jeu final utilisé}
Le jeu final utilisé pour la classification comporte 33 variables conservées après nettoyage, sélection et suppression des colonnes à risque de fuite.

% ================= PCA ET CLUSTERING =================
\section{Réduction de dimension et segmentation}

\subsection{Analyse en composantes principales (PCA)}
Une ACP automatique a été appliquée après prétraitement afin de réduire la dimensionnalité tout en conservant l’essentiel de l’information.

\begin{figure}[H]
    \centering
    \includegraphics[width=0.78\textwidth]{reports/pca_cumulative_variance.png}
    \caption{Variance cumulée expliquée par les composantes principales.}
\end{figure}

Le seuil de variance cumulée retenu est de 85\%. L’exécution finale montre qu’environ 22 composantes sont nécessaires pour atteindre ce seuil.

\subsection{Projection 2D}
\begin{figure}[H]
    \centering
    \includegraphics[width=0.78\textwidth]{reports/pca_projection_2d.png}
    \caption{Projection 2D des observations après PCA.}
\end{figure}

La projection PCA en deux dimensions montre une forte concentration centrale et quelques points plus éloignés. Cela suggère la présence d’un comportement moyen dominant et de quelques profils atypiques.

\subsection{Segmentation KMeans}
Après réduction de dimension, l’algorithme KMeans a été appliqué afin d’explorer une segmentation non supervisée des clients.

\begin{figure}[H]
    \centering
    \includegraphics[width=0.78\textwidth]{reports/kmeans_elbow.png}
    \caption{Méthode du coude pour KMeans.}
\end{figure}

\begin{figure}[H]
    \centering
    \includegraphics[width=0.78\textwidth]{reports/kmeans_silhouette.png}
    \caption{Score de silhouette selon le nombre de clusters.}
\end{figure}

Le meilleur compromis est obtenu pour \(k=2\), ce qui signifie que la segmentation la plus stable met en évidence deux grands groupes plutôt qu’une partition fine.

\begin{figure}[H]
    \centering
    \includegraphics[width=0.82\textwidth]{reports/kmeans_clusters_pca_map.png}
    \caption{Clusters KMeans projetés sur le plan PCA.}
\end{figure}

\subsection{Interprétation}
La segmentation KMeans reste principalement exploratoire. Elle permet de distinguer un groupe de clients au comportement plus central et un autre groupe plus dispersé. Néanmoins, la séparation visuelle reste modérée, ce qui confirme que la classification supervisée demeure l’axe principal du projet.

% ================= CLASSIFICATION =================
\section{Modélisation prédictive : classification}

\subsection{Objectif}
L’objectif principal de cette partie est de prédire le churn client à partir des variables préparées.

\subsection{Modèles testés}
Les modèles suivants ont été évalués :
\begin{itemize}
    \item Logistic Regression
    \item Logistic Regression avec GridSearchCV
    \item Logistic Regression avec Optuna
    \item SVM
    \item SVM avec GridSearchCV
    \item Random Forest baseline
    \item Random Forest avec GridSearchCV
\end{itemize}

\subsection{Métriques d’évaluation}
Les métriques retenues sont :
\begin{itemize}
    \item Accuracy
    \item F1-score
    \item ROC-AUC
    \item Matrice de confusion
\end{itemize}

\subsection{Résultats de classification}
Le meilleur modèle obtenu est \texttt{RandomForest\_GridSearch}.

\begin{table}[H]
\centering
\caption{Synthèse des résultats de classification}
\begin{tabular}{lccc}
\toprule
\textbf{Modèle} & \textbf{Accuracy} & \textbf{F1-score} & \textbf{ROC-AUC} \\
\midrule
LogisticRegression & 0.8697 & 0.8202 & 0.9491 \\
LogisticRegression\_GridSearch & 0.8697 & 0.8202 & 0.9491 \\
LogisticRegression\_Optuna & 0.8686 & 0.8189 & 0.9491 \\
SVM & 0.8754 & 0.8222 & 0.9510 \\
SVM\_GridSearch & 0.8914 & 0.8435 & 0.9585 \\
RandomForest\_baseline & 0.8994 & 0.8483 & 0.9602 \\
RandomForest\_GridSearch & \textbf{0.9051} & \textbf{0.8566} & \textbf{0.9614} \\
\bottomrule
\end{tabular}
\end{table}

\subsection{Matrice de confusion}
\begin{figure}[H]
    \centering
    \includegraphics[width=0.68\textwidth]{reports/confusion_matrix_best.png}
    \caption{Matrice de confusion du meilleur modèle de classification.}
\end{figure}

La matrice de confusion montre :
\begin{itemize}
    \item 544 vrais négatifs ;
    \item 248 vrais positifs ;
    \item 40 faux positifs ;
    \item 43 faux négatifs.
\end{itemize}

Le modèle classe donc correctement la grande majorité des clients fidèles et des churners.

\subsection{Interprétation}
Les performances obtenues sont élevées mais réalistes. Le modèle final est plus robuste que les premières versions quasi parfaites, ce qui confirme que la phase de détection de fuite d’information et de sélection des variables a été déterminante.

\subsection{Importance des variables}
\begin{figure}[H]
    \centering
    \includegraphics[width=0.92\textwidth]{reports/feature_importance_top20.png}
    \caption{Top 20 des variables les plus importantes pour la prédiction du churn.}
\end{figure}

L’analyse d’importance montre que le modèle repose principalement sur :
\begin{itemize}
    \item la saisonnalité (\texttt{FavoriteSeason}, \texttt{PreferredMonth}) ;
    \item l’activité transactionnelle (\texttt{Frequency}, \texttt{MonetaryTotal}) ;
    \item la diversité produit (\texttt{UniqueProducts}) ;
    \item le rythme d’achat (\texttt{AvgDaysBetweenPurchases}) ;
    \item certaines composantes temporelles d’inscription.
\end{itemize}

% ================= REGRESSION =================
\section{Modélisation prédictive : régression}

\subsection{Objectif}
L’objectif secondaire du projet est de prédire \texttt{MonetaryTotal}.

\subsection{Modèles testés}
Plusieurs modèles ont été comparés :
\begin{itemize}
    \item Ridge (baseline linéaire)
    \item RandomForestRegressor
    \item ExtraTreesRegressor
    \item HistGradientBoostingRegressor
\end{itemize}

\subsection{Métriques}
Les métriques retenues sont :
\begin{itemize}
    \item RMSE
    \item MAE
    \item \(R^2\)
\end{itemize}

\subsection{Résultats}
Le meilleur modèle obtenu est \texttt{HistGradientBoostingRegressor}.

\begin{table}[H]
\centering
\caption{Résultats de régression sur l’ensemble de test}
\begin{tabular}{lccc}
\toprule
\textbf{Modèle} & \textbf{RMSE} & \textbf{MAE} & \textbf{\(R^2\)} \\
\midrule
Ridge & 104233183.86 & 3524695.50 & -102722685.4452 \\
RandomForestRegressor & 7535.86 & 568.82 & 0.4631 \\
ExtraTreesRegressor & 7913.41 & 610.46 & 0.4079 \\
HistGradientBoostingRegressor & \textbf{6124.65} & \textbf{458.13} & \textbf{0.6453} \\
\bottomrule
\end{tabular}
\end{table}

\subsection{Analyse critique}
Le modèle Ridge a présenté un \(R^2\) fortement négatif, ce qui montre que l’hypothèse linéaire n’est pas adaptée à cette tâche. Les modèles non linéaires basés sur les arbres se sont révélés beaucoup plus performants. Le meilleur compromis est atteint avec \texttt{HistGradientBoostingRegressor}, qui capture mieux les relations complexes présentes dans les données.

% ================= VALIDATION CROISEE =================
\section{Validation croisée et robustesse}

\subsection{Validation croisée en classification}
Une validation croisée stratifiée en 5 plis a été réalisée pour les principaux modèles de classification.

\begin{table}[H]
\centering
\caption{Validation croisée en classification}
\begin{tabular}{lccc}
\toprule
\textbf{Modèle} & \textbf{Accuracy} & \textbf{F1-score} & \textbf{ROC-AUC} \\
\midrule
LogisticRegression & 0.8952 $\pm$ 0.0123 & 0.8985 $\pm$ 0.0118 & 0.9652 $\pm$ 0.0060 \\
SVM & 0.9062 $\pm$ 0.0119 & 0.9078 $\pm$ 0.0108 & 0.9716 $\pm$ 0.0071 \\
RandomForest & \textbf{0.9207 $\pm$ 0.0058} & \textbf{0.9214 $\pm$ 0.0059} & \textbf{0.9797 $\pm$ 0.0040} \\
\bottomrule
\end{tabular}
\end{table}

\subsection{Validation croisée en régression}
Une validation croisée en 5 plis a également été réalisée pour la régression.

\begin{table}[H]
\centering
\caption{Validation croisée en régression}
\begin{tabular}{lcc}
\toprule
\textbf{Modèle} & \textbf{\(R^2\)} & \textbf{RMSE} \\
\midrule
RandomForestRegressor & 0.5992 $\pm$ 0.0960 & 4078.34 $\pm$ 1820.85 \\
ExtraTreesRegressor & 0.6042 $\pm$ 0.1496 & 4107.01 $\pm$ 2035.26 \\
HistGradientBoostingRegressor & \textbf{0.7667 $\pm$ 0.1201} & \textbf{3167.55 $\pm$ 1790.01} \\
\bottomrule
\end{tabular}
\end{table}

\subsection{Interprétation}
Les résultats de validation croisée confirment la robustesse de la classification et l’amélioration nette de la régression. HistGradientBoostingRegressor apparaît comme le modèle le plus performant et le plus cohérent pour la tâche de régression.

% ================= INFERENCE =================
\section{Prédiction et inférence}

\subsection{Rôle du script de prédiction}
Le script de prédiction recharge :
\begin{itemize}
    \item le préprocesseur sauvegardé ;
    \item le meilleur modèle de classification ;
    \item les bornes IQR des outliers ;
    \item les informations de colonnes supprimées.
\end{itemize}

\subsection{Pipeline d’inférence}
Le pipeline d’inférence reproduit strictement les étapes apprises pendant l’entraînement :
\begin{itemize}
    \item parsing des dates ;
    \item feature engineering sur l’adresse IP ;
    \item suppression des colonnes inutiles ou fuyantes ;
    \item application des bornes IQR ;
    \item alignement des colonnes avec celles attendues par le modèle ;
    \item prédiction finale et estimation de probabilité.
\end{itemize}

\subsection{Résultats observés}
Les tests sur profils réalistes montrent que le modèle distingue correctement les clients stables des clients à risque, même si certaines probabilités restent proches du seuil de décision, ce qui traduit une incertitude réaliste dans des cas intermédiaires.

% ================= DEPLOIEMENT =================
\section{Déploiement avec Flask}

\subsection{Objectif}
Le déploiement vise à rendre le modèle accessible via une interface simple permettant à un utilisateur de tester différents profils clients.

\subsection{Architecture}
La solution repose sur :
\begin{itemize}
    \item une interface HTML/CSS ;
    \item un backend Flask ;
    \item le préprocesseur sauvegardé ;
    \item le meilleur modèle entraîné ;
    \item les artefacts de prétraitement.
\end{itemize}

\subsection{Modes d’utilisation}
L’interface propose :
\begin{itemize}
    \item un mode manuel à partir de 10 variables ;
    \item un mode test exact à partir d’une ligne réelle de \texttt{X\_test.csv}.
\end{itemize}

\subsection{Problèmes rencontrés}
Plusieurs difficultés ont été rencontrées au cours du projet :
\begin{itemize}
    \item fuite d’information dans certaines variables ;
    \item présence de colonnes trop redondantes ;
    \item déséquilibre de classes ;
    \item difficulté d’interprétation des modèles linéaires en régression ;
    \item écart entre le modèle complet et les entrées simplifiées de l’interface Flask ;
    \item erreurs de conversion des nombres saisis manuellement.
\end{itemize}

\subsection{Corrections apportées}
Ces problèmes ont été corrigés en :
\begin{itemize}
    \item supprimant les variables suspectes ;
    \item rééquilibrant les classes avec SMOTENC ;
    \item choisissant des modèles non linéaires plus adaptés ;
    \item alignant strictement le pipeline d’inférence sur le pipeline d’entraînement ;
    \item testant l’interface sur des profils réels et manuels.
\end{itemize}

\subsection{Limites}
L’interface actuelle reste volontairement simplifiée. En mode manuel, seules 10 variables sont renseignées ; le reste doit être reconstruit ou imputé, ce qui peut introduire un léger décalage par rapport aux prédictions obtenues sur des lignes complètes.

% ================= DISCUSSION =================
\section{Discussion générale}

\subsection{Points forts du projet}
Les principaux points forts sont :
\begin{itemize}
    \item une chaîne ML complète de bout en bout ;
    \item une gestion méthodologiquement correcte du data leakage ;
    \item un preprocessing rigoureux ;
    \item une comparaison de plusieurs modèles ;
    \item une optimisation par GridSearchCV et Optuna ;
    \item une validation croisée ;
    \item une phase d’inférence et une interface Flask.
\end{itemize}

\subsection{Limites}
Les limites principales du projet sont :
\begin{itemize}
    \item le caractère essentiellement exploratoire de la segmentation KMeans ;
    \item une séparation probabiliste parfois modérée dans les cas limites ;
    \item un décalage possible entre le modèle complet et l’interface manuelle simplifiée.
\end{itemize}

\subsection{Améliorations possibles}
Plusieurs pistes d’amélioration peuvent être envisagées :
\begin{itemize}
    \item enrichir l’interface avec davantage de variables ;
    \item calibrer les probabilités ;
    \item intégrer des outils d’explicabilité comme SHAP ;
    \item tester des modèles supplémentaires pour la régression ;
    \item développer un modèle dédié au mode manuel Flask.
\end{itemize}

% ================= CONCLUSION =================
\section{Conclusion générale}

Ce projet a permis de mettre en œuvre une démarche complète de Machine Learning appliquée à la prédiction du churn client dans un contexte retail. À partir d’un jeu de données riche mais imparfait, un pipeline robuste a été construit pour préparer les données, réduire la dimension, segmenter les clients, entraîner plusieurs modèles supervisés, comparer leurs performances et déployer une solution d’inférence.

Le meilleur modèle de classification obtenu est un \texttt{RandomForest\_GridSearch}, qui atteint de très bonnes performances tout en restant cohérent après suppression des variables à risque de fuite. En parallèle, la partie régression a été significativement améliorée grâce à \texttt{HistGradientBoostingRegressor}, mieux adapté à la structure non linéaire des données. Enfin, une interface Flask a été développée pour démontrer la faisabilité d’un usage applicatif.

L’ensemble du travail montre qu’une préparation rigoureuse des données, une méthodologie soignée et une évaluation critique sont les fondements d’un projet ML crédible, professionnel et défendable.

\end{document}